% New English translation from checked cote 25, site 12 (PDF page 13).
% AI-translated by GPT-6 Astra (OpenAI, ChatGPT Pro mode).
\scanpage{25/12}
\mspage{\hand{9}}

\uline{Corollary} 2.11. Let $y\in Y_{\hand{\eta}}$; for $Y_{\hand{\eta}}$ to \var{\inserted{be regular (resp.}}{be regular (resp.}{insertion}{ega25_25_012-v0001} satisfy property
$R_k$ (ref) at y\hand{,} \var{\deleted{(}}{}{deletion}{ega25_25_012-v0002}resp. be normal at y), it is n and s that X satisfy
the same property at x. \var{\xout{\unsure{Remark}}}{}{typewriter strike-out}{ega25_25_012-v0003} Let Z be the set of points of X at
which X is not \var{\deleted{\var{\xout{\illegible{1}}}{}{typewriter strike-out}{ega25_25_012-v0005} (}}{}{deletion}{ega25_25_012-v0004}\var{\inserted{regular (resp. \var{\deleted{$R_k$}}{}{deletion}{ega25_25_012-v0007} $R_k$}}{regular (resp. \var{\deleted{$R_k$}}{}{deletion}{ega25_25_012-v0007} $R_k$}{insertion}{ega25_25_012-v0006}resp. \var{\deleted{is not}}{}{deletion}{ega25_25_012-v0008} normal); for $Y_{\hand{\eta}}$ to be
\var{\inserted{regular (resp.}}{regular (resp.}{insertion}{ega25_25_012-v0009} $(R_k)$\hand{,} \var{\deleted{(}}{}{deletion}{ega25_25_012-v0010}resp. normal), it is n and s that f(Z) be finite i.e. dim $\overline{f(Z)}=0$.

\var{\margin{|}}{\placedmargin{|}}{marginal note placed}{ega25_25_012-v0011} The same proof as for 2.8. and 2.9. The various
references ensuring that Z is closed must be given (since we must know that it is
constructible to apply 2.3.).

\var{\xout{\unsure{Remark 2.12}}}{}{typewriter strike-out}{ega25_25_012-v0012} Observe that in 2.10. and 2.11. we say nothing
about the corresponding ``geometric'' properties, the results
considered being of an \hand{``}absolute\hand{''} nature. We now Examine
properties of a geometric nature. (One could, if desired, take the
\var{\margin{||}}{\placedmargin{||}}{marginal note placed}{ega25_25_012-v0013}opportunity to begin a new n$^{\text{o}}$).

\var{\inserted[typed]{Assume that $f:X\hand{\to} P$ is an \var{\deleted{immersion}}{}{deletion}{ega25_25_012-v0015}\var{\inserted{\unsure{unramified}}}{\unsure{unramified}}{insertion}{ega25_25_012-v0016}.}}{Assume that $f:X\hand{\to} P$ is an \var{\deleted{immersion}}{}{deletion}{ega25_25_012-v0015}\var{\inserted{\unsure{unramified}}}{\unsure{unramified}}{insertion}{ega25_25_012-v0016}.}{insertion}{ega25_25_012-v0014}
\uline{Theorem} 2.12. Let $y\in Y_{\hand{\eta}}$, x its image in X. For X to be
smooth over k at x, it is n and s that $Y_{\hand{\eta}}$ be smooth over $k(\hand{\eta})$ at y.

We may assume k algebraically closed. When $Y_{\hand{\eta}}$ is smooth over $k(\hand{\eta})$
at y, it is regular, so, since $Y_{\hand{\eta}}$ is flat over X, X is regular
at x (ref), and hence smooth over k at x since k is alg closed and therefore
perfect (ref). For the converse, we may (replacing X by
an open neighborhood of x) assume X smooth, and (by the Jacobian criterion
for smoothness) defined in an open subset U of P by p equations
as $X=V(f_1,..,f_p)$, where the differentials $df_1,..,df_p$ are
everywhere linearly independent. Introducing affine coordinates
$S_1,..,S_n$ in $\overset{\hand{\vee}}{P}$ and affine coordinates $T_1,..,T_n$ in a neighborhood
of x (by choosing a hyperplane $H_\infty$ not containing x), $Y\hand{\hookrightarrow} U_{k(\hand{\eta})}$ is
then equal to $V(f_1,..,f_p,\hand{\sum_i}S_iT_i-1\ )$, and it suffices to verify that
the differentials (relative to $k(\hand{\eta})$) of $f_1,..,f_p$, $\hand{(\sum}S_iT_i\hand{)}-1$
\tnote{25-012-1}{The inserted hypothesis retains une after immersion is struck and replaced by the doubtful unramifié/non ramifié. Both candidates mean unramified; the English keeps the remaining incomplete construction an unramified. The source's capital E in Examinons is reflected in Examine. Here and elsewhere ref is a request for a reference, not a supplied citation.}

\mspage{\hand{12}}
